\documentclass{article}
\usepackage[T2A]{fontenc}
\usepackage[utf8]{inputenc}   
\usepackage[english, russian]{babel}

% Set page size and margins
% Replace `letterpaper' with`a4paper' for UK/EU standard size
\usepackage[a4paper,top=2cm,bottom=2cm,left=2cm,right=2cm,marginparwidth=1.75cm]{geometry}

\usepackage{amsmath}
\usepackage{graphicx}
\usepackage[colorlinks=true, allcolors=blue]{hyperref}
\usepackage{amsfonts}
\usepackage{amssymb}
% \usepackage[left=1cm,right=1cm,top=1cm,bottom=1cm]{geometry}
\usepackage{hyperref}
\usepackage{seqsplit}
\usepackage[dvipsnames]{xcolor}
\usepackage{enumitem}
\usepackage{algorithm}
\usepackage{algpseudocode}
\usepackage{algorithmicx}
\usepackage{mathalfa}
\usepackage{mathrsfs}
\usepackage{dsfont}
\usepackage{caption,subcaption}
\usepackage{wrapfig}
\usepackage[stable]{footmisc}
\usepackage{indentfirst}
\usepackage{rotating}
\usepackage{pdflscape}
\usepackage{listings}

\usepackage{MnSymbol,wasysym}


\lstset{
    language=Python,
    basicstyle=\ttfamily\small,
    keywordstyle=\color{blue},
    stringstyle=\color{red},
    commentstyle=\color{gray},
    backgroundcolor=\color{lightgray!20},
    frame=single,
    rulecolor=\color{black},
    numbers=left,
    numberstyle=\tiny\color{gray},
    stepnumber=1,
    numbersep=10pt,
    showspaces=false,
    showstringspaces=false,
    breaklines=true,
    breakatwhitespace=true,
    tabsize=4,
    captionpos=b
}

\begin{document}

\begin{center}
	\Large {AI Masters \\
		Алгоритмы-2024 зима \\
		семестровая контрольная (final) \\
		Никита Загайнов}
\end{center}

\bigskip

\textbf{1[3]} Рассмотрим следующую реализацию вычисления суммы $n$ элементов. Сначала первый элемент складывается со вторым, третий с четвертым и так далее попарно. Потом этот процесс повторяется с появившимися парами и так далее, пока не останется одно число. Все числа в начале $k$-битовые, сложение линейно по битовой длине. По мере выполнения алгоритма длина чисел растет. Оцените сложность выполнения этого алгоритма.

\medskip

\textbf{Решение}

\medskip

Время алгоритма: $T(n, k) = T(\frac{n}{2}, k + 1) + nk$
Докажем, что итоговая асимптотика $O(n(k + \log n))$. Очевидно, что
максимальное значение параметра $k$ равно $k + \log_2 n$ на последнем
вызове функции. Принимая во внимание этот факт, имеем:
\[
	T(n, k) =
	O(n(k + \log n) + \frac{n}{2}(k +  \log \frac{n}{2}) +
	\ldots + (k + \log n)) = O(n(k + \log n))
\]

\medskip

\textbf{2[1 + 1]}

\begin{enumerate}
	\item $T(n) = T(n - 3) + n$
	\item $T(n) = 27 T(\frac{n}{3}) + n^4$
\end{enumerate}

\medskip

\textbf{Решение}

\medskip

\begin{enumerate}
	\item \[
		      T(n) = n + (n - 3) + (n - 6) + \ldots + 3 =
		      \sum_{i = 1}^{n / 3} 3i = O(n^2)
	      \]

	\item Применим третий случай мастер-теоремы:
	      \[
		      n^4 = \Omega (n^{\log_3 27 + \varepsilon}) = \Omega(n^{3 + \varepsilon})
	      \]
	      И
	      \[
		      27 \left(\frac{n}{3}\right)^4 = 27 \cdot \frac{n^4}{81} = \frac{1}{3} n^4 < cn^4, \quad c > \frac{1}{3}
	      \]

	      Тогда по третьему случаю мастер-теоремы:
	      \[
		      T(n) = \Theta(n^4)
	      \]
\end{enumerate}

\medskip

\textbf{3[1]} Оцените асимптотику числа печатей.

\begin{lstlisting}
def f(n):
    if (n < 100):
        return
    
    f(n / 2); f(n / 2); f(n / 2)

    for i = 0; i < n; i += 1:
        for m = 1; m < i**2; m *= 2:
            print("hehe")
\end{lstlisting}

\medskip

\textbf{Решение}

\medskip

Количество печатей:
\[
	T(n) = \begin{cases}
		0,                    & n < 100    \\
		3T(n / 2) + n \log n, & n \geq 100
	\end{cases}
\]

Очевидно, что $T(n) = O(\bar{T}(n))$ для
\[
	\bar{T}(n) = 3\bar{T}(n / 2) + n \log n
\]
По первому случаю мастер-теоремы:
\[
	n \log n = O(n^{\log_2 3 - \varepsilon}), \quad \varepsilon < \log_2 3 - 1
\]
Тогда асимптотика будет:

\[
	T(n) = \bar{T}(n) = O(n^{\log_2 3})
\]

\medskip

\textbf{4[2]} Алхимики работают с $n$ веществами, превращая их друг в друга. Вам дают алхимическую книгу, в которой описаны эти превращения. Предложите алгоритм, позволяющий найти все такие вещества, которые можно получить из вещества $A$ и такие, что из каждого из них можно получить $A$.

\medskip

\textbf{Решение}

\medskip
Превращения в книге представимы в виде графа, где вершины -- вещества,
а ребра -- превращения. Тогда чтобы найти все необходимые вещества,
нужно запустить \texttt{DFS} из вершины $A$ в исходном графе и
в обратном графе. В результате комбинируем те вершины, которые были посещены,
тогда получим искомое множество. Этот алгоритм работает, так как
первый обход найдёт все вещества, которые могут получиться из $A$
(по построению \texttt{DFS}),
а второй обход найдёт все вещества, из которых можно получить $A$
(по построению обратного графа).

Сложность -- $O(|V| + |E|)$, $|V|$ - количество веществ, $|E|$ - количество превращений.
(\texttt{DFS} - $O(|V| + |E|)$, Обращение графа - $O(|E|)$).

\medskip

\textbf{5[3]} Вы с коллегой решили найти кратчайшие пути из вершины $s$ во все остальные в ориентированном ациклическом графе с положительными целыми весами, ограниченными сверху константой $W$, в котором степени всех вершин не менее $3$. Коллега разделил все ребра на единичные и массу промежуточных вершин, не нашел пути, ушел обедать, а потом домой. Предложите алгоритм, который вернет граф в первоначальное состояние, докажите его корректность и оцените асимптотику.

\medskip

\textbf{Решение}

\medskip

Так как в исходном графе степени всех вершин не менее $3$,
то нам нужно убрать все вершины степени $2$, они были добавлены искусственно.
Это гарантированно вернёт граф в исходное состояние, так как при введении
промежуточных вершин, мы добавляем только вершины степени $2$. Сделать
это можно следующей модификацией \texttt{DFS}: хранить последнюю
рассмотренную вершину, у которой была степень $3$, и соединять её
с настоящей вершиной, если её степень тоже $3$. Далее удалить все вершины
степени $2$. Мы не создадим новых рёбер, так как если две вершины
степени $3$ соединены ребром, алгоритм добавит уже существующее ребро
(Его за новое не считаем), в обратном случае, ребро будет добавлено
между парой вершин степени $3$, увеличивая их степень и не влияя на степени
промежуточных вершин.

Сложность -- $O(W|V| + W|E|)$ -- сложность обхода в графе,
где после введения искусственных вершин добаывилось не болле $W|V|$ вершин и
$W|E|$ рёбер.

\medskip

\textbf{6[4]} Предложите алгоритм, принимающий на вход матрицу смежности неориентированного графа с единичными весами и вычисляющий матрицу весов кратчайших путей для всех пар вершин.

\medskip

\textbf{Решение}

\medskip

Воспользуемся алгоритмом Флойда-Варшалла, который высчитывает
длины кратчайших путей путём обновления исходной матрицы смежности
$d_{ij}$:
\[
	d_{ij} = \min(d_{ij}, d_{ik} + d_{kj})
\]
для всех $i, j, k$. Алгоритм корректен, так как он обновляет
расстояние между каждой парой вершин $i, j$ через промежуточную
вершину $k$ для каждой вершины $k$.

Сложность -- $O(|V|^3)$, где $|V|$ -- количество вершин в графе.

\medskip

\textbf{7[4]} На вход задачи поступает набор $\{(x_i, y_i, z_i, r_i)\}_{i=1}^n$. Каждая четверка чисел задает шар с центром в $x_i, y_i, z_i$ радиуса $r_i$. Постройте алгоритм, проверяющий, можно ли пройти из точки $x_1, y_1, z_1$ в точку $x_2, y_2, z_2$, в каждый момент времени оставаясь внутри шаров. Докажите его корректность и оцените асимптотику.

\medskip

\textbf{Решение}

\medskip

Для решения задачи построим граф, где вершины -- центры шаров, рёбра соединяют
только те шары, которые пересекаются (расстояние между центрами не больше суммы радиусов).
Далее запустим \texttt{DFS} из точки $x_1, y_1, z_1$ и проверим, достигаема ли
точка $x_2, y_2, z_2$. Алгоритм корректен по построению \texttt{DFS}. Чтобы построить
граф, потребуется проверить расстояние между каждой парой шаров, что
займёт $O(|V|^2)$ времени. \texttt{DFS} работает за $O(|V| + |E|)$.

Сложность алгоритма -- $O(|V|^2 + |E|)$.

\medskip

\textbf{8[4]} Известно, что в ориентированном графе с целыми весами есть ровно один несамопересекающийся отрицательный цикл, который при этом является стоком. Предложите алгоритм, который получает на вход описание графа и выводит число, которое нужно прибавить к весам всех ребер, чтобы убрать отрицательный цикл. Докажите его корректность и оцените асимптотику.

\medskip

\textbf{Решение}

\medskip

Так как цикл один и он является стоком, его можно найти обходом графа
(если алгоритм повторно посетил какую-то из вершин). Далее, обходим
найденный цикл, считая его вес, и возвращаем $\max (-w, 0)$. Алгоритм
корректен, так как цикл несамопересекающийся и у него только один вариант
обхода.

Сложность алгоритма -- $O(|V| + |E|)$ (два обхода графа).

\medskip

\textbf{9[5]} Рыбак взял с собой лодку, но забыл весла. Выйти из положения он намеревается с помощью прочного спиннинга, на катушке которого намотано $l$ метров прочной лески. Дело происходит в Карелии, так что на озере много ($n$) островов. Рыбак может бросить блесну на соседний остров, если он не дальше $l$ метров, и подтянуться туда путем сматывания удочки. Ему нужно добраться до конкретного острова, поскольку там лучше всего клюет. На одном из островов лежит вторая катушка длины $L > l$. Постройте алгоритм, позволяющий найти кратчайший путь до целевого острова, докажите его корректность и оцените асимптотику. Карта озера известна, остров с катушкой известен. Начать передвижение можно в любой точке берега, но после этого ступать на берег нельзя, поскольку там ходит рыбнадзор и проверяет лицензии.

\medskip

\textbf{Решение}

\medskip

Для начала построим два графа и одно множество. В обоих графах
вершинами будут острова, которые соединены рёбрами, если расстояние
между ними не превышает $l$ или $L$ (для первого и второго графа).
В множество положим все все острова, до которых можно добраться от берега.
Далее запустим \texttt{DFS} из целевой вершины в первом графе. Если
досягаема хотя бы одна вершина из множества, то рыбак может добраться
без новой лески. Если нет, то запускаем \texttt{DFS} во втором графе,
проверяя досягаемость острова с новой леской, а потом запускаем его
в первом графе из этого острова с леской, проверяя, можно ли добраться
до одного из островов из множества. Если оба раза получилось, то целевой
остров досягаем, если нет, то нет. Алгоритм корректен по построению
\texttt{DFS}.

Сложность алгоритма -- $O(|V| + |E|)$, так как мы запускаем обход не
более 3 раз.

\end{document}